\section{Aprendizaje práctico de ciberseguridad}

El aprendizaje práctico de ciberseguridad surge de la necesidad de que el estudiante no solo comprenda conceptos, sino que pueda aplicarlos frente a situaciones concretas. Después de reconocer principios como la confidencialidad, la integridad y la disponibilidad, así como amenazas y vulnerabilidades frecuentes, resulta necesario analizar cómo esos contenidos pueden trabajarse desde una perspectiva formativa. Por ello, este capítulo aborda las limitaciones del enfoque exclusivamente teórico, el aprendizaje activo, el uso de casos, la práctica guiada y la retroalimentación como elementos relevantes para fortalecer habilidades aplicadas en ciberseguridad.

\subsection{Limitaciones del enfoque teórico}

La enseñanza de la ciberseguridad puede apoyarse en clases expositivas, lecturas y evaluaciones conceptuales. Estos recursos son importantes porque permiten construir una base inicial de conocimiento; sin embargo, resultan limitados cuando el propósito es desarrollar habilidades aplicadas. Reconocer un intento de \textit{phishing}, interpretar una señal de riesgo o responder ante una situación sospechosa exige análisis, práctica y toma de decisiones, no únicamente memorización de definiciones.

En áreas vinculadas a ciencia, tecnología, ingeniería y matemáticas, las estrategias de aprendizaje activo han mostrado mejores resultados que la enseñanza centrada únicamente en exposición. Freeman et al. (2014) señalan que el aprendizaje activo incrementa el rendimiento académico y reduce las tasas de reprobación en comparación con clases tradicionales, debido a que involucra al estudiante en actividades, discusión y resolución de problemas. Esta evidencia permite sostener que, en un área aplicada como la ciberseguridad, la teoría debe complementarse con actividades donde el estudiante participe de manera directa.

En ciberseguridad, la distancia entre teoría y práctica puede hacerse evidente cuando el estudiante conoce un concepto, pero no logra aplicarlo frente a un caso concreto. Por ejemplo, puede saber qué es una amenaza digital, pero no necesariamente identificarla en una situación cotidiana. Por esta razón, el proceso formativo debe incorporar actividades que permitan observar, comparar, decidir y corregir.

\subsection{Aprendizaje activo}

El aprendizaje activo plantea que el estudiante participa directamente en la construcción de su conocimiento mediante actividades, análisis de problemas, discusión de casos y toma de decisiones. A diferencia de un enfoque pasivo, donde el estudiante solo recibe información, este enfoque promueve que el aprendizaje se produzca a través de la participación y la aplicación de conceptos en tareas concretas (Freeman et al., 2014).

En el caso de la ciberseguridad, este enfoque resulta pertinente porque los contenidos están relacionados con situaciones variables y con riesgos que pueden depender del contexto. Un mismo concepto, como el uso indebido de credenciales o la ingeniería social, puede presentarse de distintas formas. Por ello, el estudiante necesita practicar la interpretación de señales, la evaluación de alternativas y la elección de respuestas adecuadas.

La práctica activa favorece una comprensión más profunda que la simple repetición de definiciones. Cuando el estudiante analiza una situación, identifica elementos de riesgo y elige una acción, desarrolla criterios que pueden ser reutilizados en escenarios similares. Esta forma de aprendizaje permite vincular los conceptos de seguridad con decisiones concretas.

\subsection{Aprendizaje basado en casos}

El aprendizaje basado en casos permite trabajar con situaciones específicas que representan problemas, decisiones o eventos relacionados con un área de conocimiento. En ciberseguridad, este enfoque es especialmente útil porque muchas amenazas no aparecen como conceptos abstractos, sino como situaciones concretas: un enlace sospechoso, una contraseña débil, un archivo desconocido, una solicitud engañosa o una configuración insegura.

Los casos permiten conectar los fundamentos de la ciberseguridad con problemas reconocibles para el estudiante. Al analizar un caso, no solo se identifica una amenaza o vulnerabilidad, sino que también se revisan sus posibles consecuencias. De esta manera, el estudiante puede comprender por qué una acción aparentemente simple, como abrir un enlace o compartir una credencial, puede comprometer la seguridad de un sistema.

En la formación de ciberseguridad, los escenarios prácticos son relevantes porque permiten representar situaciones de riesgo sin exponer sistemas reales. Vykopal et al. (2021) señalan que los entornos prácticos de entrenamiento permiten ejercitar conceptos y herramientas en condiciones más realistas, aunque también advierten que su preparación y aplicación puede ser compleja cuando se trabaja con grupos numerosos. Por ello, el uso de casos organizados y reutilizables puede apoyar una enseñanza más clara y progresiva.

\subsection{Práctica guiada}

La práctica guiada consiste en desarrollar actividades aplicadas con algún tipo de orientación, ya sea mediante instrucciones, ejemplos, pistas, retroalimentación o revisión posterior. En ciberseguridad, esta orientación es importante porque el estudiante puede cometer errores durante el aprendizaje, pero esos errores deben ocurrir en un ambiente controlado y sin afectar sistemas reales.

Una actividad práctica sin orientación suficiente puede generar confusión, especialmente cuando el estudiante todavía no domina los conceptos básicos. Por esta razón, la guía durante el proceso permite ordenar la experiencia de aprendizaje y reducir la posibilidad de que el estudiante solo actúe por ensayo y error. En este sentido, la práctica guiada ayuda a vincular los conceptos previamente estudiados con decisiones más razonadas.

La enseñanza práctica de ciberseguridad también presenta dificultades de escala. Vykopal et al. (2021) explican que preparar clases prácticas puede requerir esfuerzo considerable, sobre todo cuando se necesitan ambientes de trabajo, tareas específicas, retroalimentación y análisis del aprendizaje. Este aspecto es relevante en contextos universitarios con grupos grandes, donde el seguimiento individual puede ser limitado.

\subsection{Retroalimentación en el aprendizaje}

La retroalimentación es un componente central del aprendizaje porque permite que el estudiante comprenda el efecto de sus acciones y corrija errores durante el proceso. Hattie y Timperley (2007) explican que una retroalimentación efectiva ayuda a responder tres preguntas: hacia dónde se dirige el estudiante, cómo está avanzando y qué necesita hacer después. En ciberseguridad, estas preguntas son útiles porque muchas decisiones requieren ajuste, revisión y mejora.

En una actividad práctica, la retroalimentación no debería limitarse a indicar si una respuesta es correcta o incorrecta. También puede explicar por qué una decisión fue riesgosa, qué señal debió observarse o qué alternativa habría sido más adecuada. De esta manera, el estudiante no solo recibe una calificación, sino una orientación para mejorar su criterio ante situaciones similares.

La retroalimentación inmediata resulta especialmente valiosa cuando el estudiante enfrenta ejercicios vinculados con amenazas digitales, porque permite corregir errores antes de que se consoliden malas prácticas. Si el estudiante recibe una explicación poco después de actuar, puede relacionar de manera más clara su decisión con la consecuencia generada.
